% LaTeX Template for AI Problem Analysis Output
% AMC12B 2024 Problem 21 - Strategic Analysis
% Framework Version: 2.1 (Enhanced for COMC)

\documentclass[12pt]{article}
\usepackage[utf8]{inputenc}
\usepackage{amsmath, amssymb, amsthm}
\usepackage{geometry}
\usepackage{xcolor}
\usepackage[most]{tcolorbox}
\usepackage{enumitem}
\usepackage{fancyhdr}

% Page setup
\geometry{margin=1in}
\setlength{\headheight}{14.5pt}
\pagestyle{fancy}
\fancyhf{}
\rhead{AMC12 Problem Analysis}
\lhead{2024 AMC 12B Problem 21}
\cfoot{\thepage}

% Color definitions
\definecolor{problemconfig}{RGB}{230, 242, 255}    % Light blue
\definecolor{strategic}{RGB}{255, 230, 242}       % Light pink
\definecolor{tactical}{RGB}{242, 255, 230}        % Light green
\definecolor{techniques}{RGB}{255, 242, 230}      % Light yellow
\definecolor{verification}{RGB}{255, 230, 230}    % Light red
\definecolor{solution}{RGB}{230, 255, 255}        % Light cyan
\definecolor{competition}{RGB}{242, 242, 255}     % Light purple
\definecolor{gold}{RGB}{218, 165, 32}

% Box styles
\newtcolorbox{problemconfigbox}{
    colback=problemconfig,
    colframe=blue,
    title=PROBLEM CONFIGURATION ANALYSIS,
    fonttitle=\bfseries,
    arc=3pt,
    boxrule=1pt,
    breakable
}

\newtcolorbox{strategicbox}{
    colback=strategic,
    colframe=purple,
    title=STRATEGIC FRAMEWORK APPLICATION,
    fonttitle=\bfseries,
    arc=3pt,
    boxrule=1pt,
    breakable
}

\newtcolorbox{tacticalbox}{
    colback=tactical,
    colframe=green,
    title=TACTICAL METHODOLOGY SELECTION,
    fonttitle=\bfseries,
    arc=3pt,
    boxrule=1pt,
    breakable
}

\newtcolorbox{techniquesbox}{
    colback=techniques,
    colframe=orange,
    title=CONVERSION TECHNIQUES APPLICATION,
    fonttitle=\bfseries,
    arc=3pt,
    boxrule=1pt,
    breakable
}

\newtcolorbox{verificationbox}{
    colback=verification,
    colframe=red,
    title=VERIFICATION STRATEGY DESIGN,
    fonttitle=\bfseries,
    arc=3pt,
    boxrule=1pt,
    breakable
}

\newtcolorbox{solutionbox}{
    colback=solution,
    colframe=cyan,
    title=SOLUTION ROADMAP,
    fonttitle=\bfseries,
    arc=3pt,
    boxrule=1pt,
    breakable
}

\newtcolorbox{competitionbox}{
    colback=competition,
    colframe=violet,
    title=COMPETITION STRATEGY INTEGRATION,
    fonttitle=\bfseries,
    arc=3pt,
    boxrule=1pt,
    breakable
}

\newtcolorbox{keyinsightbox}{
    colback=yellow!20,
    colframe=gold,
    title=KEY INSIGHT,
    fonttitle=\bfseries,
    arc=3pt,
    boxrule=2pt,
    leftrule=5pt,
    breakable
}

\newtcolorbox{warningbox}{
    colback=red!10,
    colframe=red!50,
    title=WARNING: COMMON PITFALL,
    fonttitle=\bfseries,
    arc=3pt,
    boxrule=2pt,
    leftrule=5pt,
    breakable
}

\newtcolorbox{tipbox}{
    colback=green!10,
    colframe=green!50,
    title=PRO TIP,
    fonttitle=\bfseries,
    arc=3pt,
    boxrule=2pt,
    leftrule=5pt,
    breakable
}

\begin{document}

\title{AMC12B 2024 Problem 21\\Strategic Problem Analysis}
\author{Competition Math Framework v2.1}
\date{\today}
\maketitle

\section*{Problem Statement}

The measures of the smallest angles of three different right triangles sum to $90^{\circ}$. All three triangles have side lengths that are primitive Pythagorean triples. Two of them are $3-4-5$ and $5-12-13$. What is the perimeter of the third triangle?

$$\textbf{(A) } 40 \qquad\textbf{(B) } 126 \qquad\textbf{(C) } 154 \qquad\textbf{(D) } 176 \qquad\textbf{(E) } 208$$

\vspace{0.3cm}
\noindent\textit{Note: The answer is not revealed here to allow students to work through the problem. See the Final Answer section at the end.}

\begin{problemconfigbox}
\subsection*{A. Problem Classification}
\begin{itemize}[leftmargin=*]
    \item \textbf{Problem Type}: To Find (perimeter of triangle)
    \item \textbf{Competition Context}: AMC12B 2024, Problem 21
    \item \textbf{Difficulty Band}: Late-band (P21) - Requires angle relationships and Pythagorean triples
    \item \textbf{Mathematical Domain}: Geometry (right triangles), Trigonometry (angle addition), Number Theory (Pythagorean triples)
    \item \textbf{Estimated Time}: 4-6 minutes
    \item \textbf{Point Value}: 6 points (standard AMC12 scoring)
\end{itemize}

\subsection*{B. Problem Structure Identification}
\begin{itemize}[leftmargin=*]
    \item \textbf{Given Information}: 
    \begin{itemize}
        \item Three right triangles with primitive Pythagorean triples
        \item Smallest angles sum to $90^{\circ}$
        \item Two known triangles: $3-4-5$ and $5-12-13$
        \item Need to find perimeter of third triangle
    \end{itemize}
    
    \item \textbf{Constraints}: 
    \begin{itemize}
        \item All triangles are right triangles
        \item All have primitive Pythagorean triples (coprime integers)
        \item Sum of smallest angles is exactly $90^{\circ}$
        \item The smallest angle in a right triangle is opposite the smallest leg
    \end{itemize}
    
    \item \textbf{Objective}: Find perimeter of third triangle
    
    \item \textbf{Hidden Assumptions}: 
    \begin{itemize}
        \item Smallest angle in a right triangle is opposite the smallest leg
        \item For $3-4-5$: smallest angle $\alpha$ has $\tan\alpha = \frac{3}{4}$
        \item For $5-12-13$: smallest angle $\beta$ has $\tan\beta = \frac{5}{12}$
        \item If $\alpha + \beta + \theta = 90^{\circ}$, then $\theta = 90^{\circ} - (\alpha + \beta)$
        \item The tangent addition formula leads to finding $\tan\theta$
    \end{itemize}
    
    \item \textbf{Answer Choice Leverage}: 
    \begin{itemize}
        \item All answers are perimeters (sums of three sides)
        \item Can eliminate some by checking if they form valid Pythagorean triples
        \item Answer (A) 40 is too small (less than known triangles)
        \item Answers cluster in two groups: small (40, 126, 154) and large (176, 208)
    \end{itemize}
    
    \item \textbf{Proof Requirements}: N/A (computational problem)
\end{itemize}

\subsection*{C. Strategic Orientation}
\begin{itemize}[leftmargin=*]
    \item \textbf{Hypothesis}: Use angle addition formulas to find the third angle, then determine the Pythagorean triple
    
    \item \textbf{Conclusion}: Need perimeter (sum of three sides)
    
    \item \textbf{Notation}: 
    \begin{itemize}
        \item $\alpha$ = smallest angle in $3-4-5$ triangle
        \item $\beta$ = smallest angle in $5-12-13$ triangle
        \item $\theta$ = smallest angle in unknown triangle
        \item Constraint: $\alpha + \beta + \theta = 90^{\circ}$
    \end{itemize}
    
    \item \textbf{Intuitive Approach}: 
    \begin{enumerate}
        \item Find $\tan\alpha = \frac{3}{4}$ and $\tan\beta = \frac{5}{12}$
        \item Use tangent addition: $\tan(\alpha + \beta)$
        \item Since $\alpha + \beta + \theta = 90^{\circ}$, we have $\tan(90^{\circ}) = \tan(\alpha + \beta + \theta)$
        \item The tangent of $90^{\circ}$ is undefined, so denominator must be zero
        \item Solve for $\tan\theta$
        \item Identify the Pythagorean triple from $\tan\theta$
    \end{enumerate}
    
    \item \textbf{Cross-Domain Potential}: 
    \begin{itemize}
        \item Geometry $\rightarrow$ Trigonometry (angles and ratios)
        \item Trigonometry $\rightarrow$ Number Theory (Pythagorean triples)
        \item Alternative: Complex numbers (arguments represent angles)
    \end{itemize}
\end{itemize}
\end{problemconfigbox}

\begin{strategicbox}
\subsection*{A. Psychological Strategies}
\begin{itemize}[leftmargin=*]
    \item \textbf{Mental Toughness}: This is P21 with multiple solution paths. The tangent addition formula is standard, but recognizing the "undefined tangent" condition is key.
    
    \item \textbf{Creativity Cultivation}: Multiple elegant approaches exist: tangent addition, complex numbers, sine/cosine, similarity. Choose what you're most comfortable with.
    
    \item \textbf{Iterative Refinement}: If tangent addition seems messy, try sine or cosine addition formulas instead.
\end{itemize}

\subsection*{B. Investigation Strategies}
\begin{itemize}[leftmargin=*]
    \item \textbf{Startup Orientation}: Recognize that "smallest angles" means angles opposite smallest legs in right triangles.
    
    \item \textbf{Get Your Hands Dirty}: 
    \begin{itemize}
        \item Calculate: $\tan\alpha = \frac{3}{4}$, $\tan\beta = \frac{5}{12}$
        \item Try: $\tan(\alpha + \beta) = \frac{\tan\alpha + \tan\beta}{1 - \tan\alpha \tan\beta}$
        \item Observe: Sum equals $90^{\circ}$, so tangent should be undefined!
    \end{itemize}
    
    \item \textbf{Penultimate Step}: Once we have $\tan\theta = \frac{33}{56}$, find the Pythagorean triple, then sum
    
    \item \textbf{Make It Easier}: Use complementary angles: $\theta = 90^{\circ} - (\alpha + \beta)$, so $\sin\theta = \cos(\alpha + \beta)$
    
    \item \textbf{Conjecture via Small Cases}: Check that $\frac{33}{56}$ gives a primitive Pythagorean triple
    
    \item \textbf{Wishful Thinking}: "I wish I knew $\tan\theta$..." $\Rightarrow$ Use the fact that $\tan(90^{\circ})$ is undefined
    
    \item \textbf{Visualization}: Draw the three right triangles with their smallest angles labeled
\end{itemize}

\subsection*{C. Late-Band Strategic Playbook}
\begin{itemize}[leftmargin=*]
    \item \textbf{Pattern Recognition}: "Sum to $90^{\circ}$" suggests complementary angles
    
    \item \textbf{Answer Choice Elimination}: 
    \begin{itemize}
        \item (A) 40 is too small
        \item Once we know one leg is 33 and hypotenuse is 65, only (C) works
    \end{itemize}
    
    \item \textbf{Tangent Addition}: Standard technique for angle sum problems
\end{itemize}

\subsection*{D. Argument Construction Strategy}
\begin{itemize}[leftmargin=*]
    \item \textbf{Direct Calculation}: Use tangent addition formula to find third angle
    
    \item \textbf{Alternative Methods}:
    \begin{itemize}
        \item Sine/cosine addition formulas
        \item Complex numbers (arguments represent angles)
        \item Coordinate geometry
        \item Similarity
    \end{itemize}
\end{itemize}

\subsection*{E. Cross-Domain Translation}
\begin{itemize}[leftmargin=*]
    \item \textbf{Geometric to Trigonometric}: Translate angle relationships into trig equations
    
    \item \textbf{Trigonometric to Algebraic}: Solve for $\tan\theta$ algebraically
    
    \item \textbf{Algebraic to Number-Theoretic}: Identify Pythagorean triple from ratio
\end{itemize}
\end{strategicbox}

\begin{tacticalbox}
\subsection*{A. Core Mathematical Tactics}
\begin{itemize}[leftmargin=*]
    \item \textbf{Primary Tactics}: 
    \begin{itemize}
        \item \textbf{Tangent Addition Formula}: Key technique for angle sums
        \item \textbf{Complementary Angles}: $\theta = 90^{\circ} - (\alpha + \beta)$
        \item \textbf{Undefined Tangent}: $\tan(90^{\circ})$ is undefined
        \item \textbf{Pythagorean Triple Recognition}: From $\tan\theta$, identify the triple
    \end{itemize}
    
    \item \textbf{Domain-Specific Tactics}:
    \begin{itemize}
        \item \textbf{Right Triangle Angle Relationships}:
        \begin{itemize}
            \item Smallest angle is opposite smallest leg
            \item For right triangle with legs $a < b$ and hypotenuse $c$: $\tan(\text{smallest angle}) = \frac{a}{b}$
        \end{itemize}
        
        \item \textbf{Tangent Addition Formula}:
        $$\tan(\alpha + \beta) = \frac{\tan\alpha + \tan\beta}{1 - \tan\alpha \tan\beta}$$
        
        \item \textbf{Complementary Angles}:
        \begin{itemize}
            \item $\sin(90^{\circ} - x) = \cos x$
            \item $\cos(90^{\circ} - x) = \sin x$
            \item $\tan(90^{\circ} - x) = \cot x = \frac{1}{\tan x}$
        \end{itemize}
        
        \item \textbf{Sine/Cosine Addition}:
        \begin{align*}
            \sin(\alpha + \beta) &= \sin\alpha \cos\beta + \cos\alpha \sin\beta \\
            \cos(\alpha + \beta) &= \cos\alpha \cos\beta - \sin\alpha \sin\beta
        \end{align*}
        
        \item \textbf{Primitive Pythagorean Triples}:
        \begin{itemize}
            \item Sides are coprime (GCD = 1)
            \item Generated by Euclid's formula: $(m^2 - n^2, 2mn, m^2 + n^2)$ where $m > n$, $\gcd(m,n) = 1$
            \item Common triples: $(3,4,5)$, $(5,12,13)$, $(8,15,17)$, $(7,24,25)$, $(20,21,29)$, $(9,40,41)$, $(12,35,37)$, $(11,60,61)$, $(13,84,85)$, $(33,56,65)$
        \end{itemize}
    \end{itemize}
    
    \item \textbf{Crossover Tactics}:
    \begin{itemize}
        \item Geometry (triangles) + Trigonometry (angles)
        \item Trigonometry + Number Theory (Pythagorean triples)
    \end{itemize}
\end{itemize}
\end{tacticalbox}

\begin{techniquesbox}
\subsection*{A. High-Probability Techniques}
\begin{itemize}[leftmargin=*]
    \item \textbf{Primary Technique}: \textbf{Tangent Addition Formula + Undefined Tangent}
    
    \item \textbf{Recognition Cues}: 
    \begin{itemize}
        \item "Smallest angles sum to $90^{\circ}$" $\Rightarrow$ Use angle addition
        \item "Right triangles" $\Rightarrow$ Tangent relates legs
        \item "Pythagorean triples" $\Rightarrow$ Integer side lengths
        \item Sum equals $90^{\circ}$ $\Rightarrow$ Tangent is undefined
    \end{itemize}
    
    \item \textbf{Application}:
    \begin{enumerate}
        \item Find $\tan\alpha = \frac{3}{4}$ (smallest leg over next smallest)
        \item Find $\tan\beta = \frac{5}{12}$
        \item Calculate $\tan(\alpha + \beta) = \frac{\frac{3}{4} + \frac{5}{12}}{1 - \frac{3}{4} \cdot \frac{5}{12}} = \frac{\frac{14}{12}}{\frac{33}{48}} = \frac{56}{33}$
        \item Use $\tan(90^{\circ}) = \tan(\alpha + \beta + \theta)$
        \item Apply tangent addition again: denominator must be 0
        \item Get $1 - \frac{56}{33} \tan\theta = 0$, so $\tan\theta = \frac{33}{56}$
        \item Triangle has legs 33 and 56, hypotenuse $\sqrt{33^2 + 56^2} = 65$
        \item Perimeter: $33 + 56 + 65 = 154$
    \end{enumerate}
\end{itemize}

\subsection*{B. Alternative Techniques}
\begin{itemize}[leftmargin=*]
    \item \textbf{Sine/Cosine Addition Approach}:
    \begin{itemize}
        \item Use $\sin\theta = \sin(90^{\circ} - (\alpha + \beta)) = \cos(\alpha + \beta)$
        \item Calculate $\cos(\alpha + \beta)$ using addition formula
        \item From $\cos\theta = \frac{56}{65}$ (adjacent over hypotenuse), identify triangle
    \end{itemize}
    
    \item \textbf{Complex Numbers Approach}:
    \begin{itemize}
        \item Represent angles as arguments: $(4 + 3i)(12 + 5i)(b + ai) = ni$ for some real $n$
        \item Real part must be zero: $56a - 33b = 0$
        \item Get $a/b = 33/56$, so legs are 33 and 56
        \item Very elegant algebraically!
    \end{itemize}
    
    \item \textbf{Answer Choice Testing}:
    \begin{itemize}
        \item Once we know one leg is 33 (or 56), check which answer works
        \item (A) 40: impossible (too small)
        \item (B) 126: would give $126 - 33 = 93$ for other two sides
        \item (C) 154: gives $154 - 33 = 121$ for other two sides, check if Pythagorean
        \item This method works but is slower
    \end{itemize}
\end{itemize}
\end{techniquesbox}

\begin{keyinsightbox}
\textbf{The Critical Insights}:

\begin{enumerate}
    \item \textbf{Smallest Angle Identification}: In a right triangle, the smallest angle is opposite the smallest leg. For $3-4-5$: $\tan\alpha = \frac{3}{4}$. For $5-12-13$: $\tan\beta = \frac{5}{12}$.
    
    \item \textbf{Tangent Addition for Two Angles}:
    \begin{align*}
        \tan(\alpha + \beta) &= \frac{\tan\alpha + \tan\beta}{1 - \tan\alpha \tan\beta} \\
        &= \frac{\frac{3}{4} + \frac{5}{12}}{1 - \frac{3}{4} \cdot \frac{5}{12}} \\
        &= \frac{\frac{9 + 5}{12}}{\frac{48 - 15}{48}} \\
        &= \frac{\frac{14}{12}}{\frac{33}{48}} \\
        &= \frac{14 \cdot 48}{12 \cdot 33} \\
        &= \frac{56}{33}
    \end{align*}
    
    \item \textbf{The Undefined Tangent Condition}: Since $\alpha + \beta + \theta = 90^{\circ}$, we have:
    $$\tan(90^{\circ}) = \tan(\alpha + \beta + \theta)$$
    
    The tangent of $90^{\circ}$ is undefined. Using the tangent addition formula:
    $$\tan(\alpha + \beta + \theta) = \frac{\tan(\alpha + \beta) + \tan\theta}{1 - \tan(\alpha + \beta) \cdot \tan\theta}$$
    
    For this to be undefined, the denominator must equal zero:
    $$1 - \tan(\alpha + \beta) \cdot \tan\theta = 0$$
    $$1 - \frac{56}{33} \cdot \tan\theta = 0$$
    
    \item \textbf{Solving for the Third Angle}:
    $$\tan\theta = \frac{33}{56}$$
    
    This means the triangle has legs in ratio $33:56$.
    
    \item \textbf{Primitive Pythagorean Triple}: Since $\gcd(33, 56) = 1$, the triangle has legs 33 and 56. The hypotenuse is:
    $$c = \sqrt{33^2 + 56^2} = \sqrt{1089 + 3136} = \sqrt{4225} = 65$$
    
    We can verify: $33^2 + 56^2 = 1089 + 3136 = 4225 = 65^2$ \checkmark
    
    \item \textbf{Perimeter Calculation}:
    $$P = 33 + 56 + 65 = 154$$
    
    \item \textbf{Alternative via Complementary Angles}: Using $\theta = 90^{\circ} - (\alpha + \beta)$:
    \begin{align*}
        \cos\theta &= \cos(90^{\circ} - (\alpha + \beta)) = \sin(\alpha + \beta) \\
        &= \sin\alpha \cos\beta + \cos\alpha \sin\beta \\
        &= \frac{3}{5} \cdot \frac{12}{13} + \frac{4}{5} \cdot \frac{5}{13} \\
        &= \frac{36 + 20}{65} = \frac{56}{65}
    \end{align*}
    
    So the adjacent leg is 56 and hypotenuse is 65, giving opposite leg $= 33$.
\end{enumerate}
\end{keyinsightbox}

\begin{verificationbox}
\subsection*{A. Verification Level Selection}
\begin{itemize}[leftmargin=*]
    \item \textbf{Chosen Level}: Level 4-5 (Standard to Thorough Verification) - 1 minute
    
    \item \textbf{Justification}: 
    \begin{itemize}
        \item Problem 21 is late-band but straightforward once you know the approach
        \item Main risks: arithmetic errors in tangent addition, Pythagorean theorem calculation
        \item Should verify the Pythagorean triple and final sum
        \item Can check answer against choices
    \end{itemize}
    
    \item \textbf{Time Budget}: 1 minute for verification
\end{itemize}

\subsection*{B. Verification Checklist}
\begin{itemize}[leftmargin=*]
    \item \textbf{Tangent Calculations}:
    \begin{itemize}
        \item Verify: $\tan\alpha = \frac{3}{4}$ (smallest leg / next leg) \checkmark
        \item Verify: $\tan\beta = \frac{5}{12}$ \checkmark
    \end{itemize}
    
    \item \textbf{Tangent Addition}:
    \begin{itemize}
        \item Check numerator: $\frac{3}{4} + \frac{5}{12} = \frac{9}{12} + \frac{5}{12} = \frac{14}{12}$ \checkmark
        \item Check denominator: $1 - \frac{3}{4} \cdot \frac{5}{12} = 1 - \frac{15}{48} = \frac{33}{48}$ \checkmark
        \item Result: $\frac{14/12}{33/48} = \frac{14 \cdot 48}{12 \cdot 33} = \frac{672}{396} = \frac{56}{33}$ \checkmark
    \end{itemize}
    
    \item \textbf{Third Angle}:
    \begin{itemize}
        \item From $1 - \frac{56}{33} \tan\theta = 0$: $\tan\theta = \frac{33}{56}$ \checkmark
    \end{itemize}
    
    \item \textbf{Pythagorean Triple}:
    \begin{itemize}
        \item Check: $33^2 + 56^2 = 1089 + 3136 = 4225 = 65^2$ \checkmark
        \item Verify primitive: $\gcd(33, 56) = \gcd(33, 56)$
        \item $33 = 3 \cdot 11$, $56 = 8 \cdot 7$, so $\gcd = 1$ \checkmark
    \end{itemize}
    
    \item \textbf{Perimeter}:
    \begin{itemize}
        \item Sum: $33 + 56 + 65 = 89 + 65 = 154$ \checkmark
        \item Matches choice (C)
    \end{itemize}
\end{itemize}

\subsection*{C. Domain-Specific Verification}
\begin{itemize}[leftmargin=*]
    \item \textbf{Angle Sum Check}:
    \begin{itemize}
        \item Can verify that $\alpha + \beta + \theta = 90^{\circ}$ using calculators
        \item $\arctan(\frac{3}{4}) + \arctan(\frac{5}{12}) + \arctan(\frac{33}{56}) \approx 36.87^{\circ} + 22.62^{\circ} + 30.51^{\circ} = 90^{\circ}$ \checkmark
    \end{itemize}
    
    \item \textbf{Alternative Method Verification}:
    \begin{itemize}
        \item Using sine: $\sin\theta = \frac{33}{65}$, $\cos\theta = \frac{56}{65}$
        \item Check: $\sin^2\theta + \cos^2\theta = \frac{1089 + 3136}{4225} = \frac{4225}{4225} = 1$ \checkmark
    \end{itemize}
\end{itemize}
\end{verificationbox}

\begin{solutionbox}
\subsection*{A. Step-by-Step Solution (Tangent Addition)}

\textbf{Step 1: Identify the Smallest Angles}

In a right triangle, the smallest angle is opposite the smallest leg.

For the $3-4-5$ triangle:
$$\tan\alpha = \frac{\text{opposite}}{\text{adjacent}} = \frac{3}{4}$$

For the $5-12-13$ triangle:
$$\tan\beta = \frac{\text{opposite}}{\text{adjacent}} = \frac{5}{12}$$

Let $\theta$ be the smallest angle in the unknown triangle.

Given: $\alpha + \beta + \theta = 90^{\circ}$

\textbf{Checkpoint}: We have the tangent values for two of the three angles.

\vspace{0.3cm}

\textbf{Step 2: Apply Tangent Addition Formula}

First, find $\tan(\alpha + \beta)$ using:
$$\tan(\alpha + \beta) = \frac{\tan\alpha + \tan\beta}{1 - \tan\alpha \tan\beta}$$

Calculate the numerator:
$$\tan\alpha + \tan\beta = \frac{3}{4} + \frac{5}{12} = \frac{9}{12} + \frac{5}{12} = \frac{14}{12} = \frac{7}{6}$$

Calculate the denominator:
\begin{align*}
    1 - \tan\alpha \tan\beta &= 1 - \frac{3}{4} \cdot \frac{5}{12} \\
    &= 1 - \frac{15}{48} \\
    &= \frac{48 - 15}{48} \\
    &= \frac{33}{48} \\
    &= \frac{11}{16} \quad \text{(simplifying by dividing by 3)}
\end{align*}

Therefore:
$$\tan(\alpha + \beta) = \frac{7/6}{11/16} = \frac{7}{6} \cdot \frac{16}{11} = \frac{112}{66} = \frac{56}{33}$$

\textbf{Checkpoint}: We have $\tan(\alpha + \beta) = \frac{56}{33}$.

\vspace{0.3cm}

\textbf{Step 3: Use the Undefined Tangent Condition}

Since $\alpha + \beta + \theta = 90^{\circ}$, we have:
$$\tan(90^{\circ}) = \tan((\alpha + \beta) + \theta)$$

The tangent of $90^{\circ}$ is undefined. Using the tangent addition formula:
$$\tan((\alpha + \beta) + \theta) = \frac{\tan(\alpha + \beta) + \tan\theta}{1 - \tan(\alpha + \beta) \cdot \tan\theta}$$

For this to be undefined, the denominator must equal zero:
$$1 - \tan(\alpha + \beta) \cdot \tan\theta = 0$$

Substituting $\tan(\alpha + \beta) = \frac{56}{33}$:
$$1 - \frac{56}{33} \cdot \tan\theta = 0$$

Solving for $\tan\theta$:
$$\frac{56}{33} \cdot \tan\theta = 1$$
$$\tan\theta = \frac{33}{56}$$

\textbf{Checkpoint}: The smallest angle in the third triangle has $\tan\theta = \frac{33}{56}$.

\vspace{0.3cm}

\textbf{Step 4: Determine the Pythagorean Triple}

Since $\tan\theta = \frac{33}{56}$, the legs of the triangle are in ratio $33:56$.

Check if this is primitive: $\gcd(33, 56) = \gcd(3 \cdot 11, 8 \cdot 7) = 1$ \checkmark

So the legs are $a = 33$ and $b = 56$.

Find the hypotenuse using the Pythagorean theorem:
\begin{align*}
    c &= \sqrt{33^2 + 56^2} \\
    &= \sqrt{1089 + 3136} \\
    &= \sqrt{4225} \\
    &= 65
\end{align*}

\textbf{Verification}: $33^2 + 56^2 = 1089 + 3136 = 4225 = 65^2$ \checkmark

\textbf{Checkpoint}: The third triangle is $33-56-65$.

\vspace{0.3cm}

\textbf{Step 5: Calculate the Perimeter}

$$P = 33 + 56 + 65 = 154$$

\textbf{Final Answer}: $\boxed{\textbf{(C) } 154}$

\subsection*{B. Alternative Solution (Sine/Cosine Addition)}

Since $\alpha + \beta + \theta = 90^{\circ}$, we have $\theta = 90^{\circ} - (\alpha + \beta)$.

Using the complementary angle relationship:
$$\sin\theta = \sin(90^{\circ} - (\alpha + \beta)) = \cos(\alpha + \beta)$$

Apply the cosine addition formula:
$$\cos(\alpha + \beta) = \cos\alpha \cos\beta - \sin\alpha \sin\beta$$

For the $3-4-5$ triangle: $\sin\alpha = \frac{3}{5}$, $\cos\alpha = \frac{4}{5}$

For the $5-12-13$ triangle: $\sin\beta = \frac{5}{13}$, $\cos\beta = \frac{12}{13}$

Therefore:
\begin{align*}
    \cos(\alpha + \beta) &= \frac{4}{5} \cdot \frac{12}{13} - \frac{3}{5} \cdot \frac{5}{13} \\
    &= \frac{48}{65} - \frac{15}{65} \\
    &= \frac{33}{65}
\end{align*}

So $\sin\theta = \frac{33}{65}$, meaning in the right triangle with smallest angle $\theta$:
\begin{itemize}
    \item Opposite leg = 33
    \item Hypotenuse = 65
    \item Adjacent leg = $\sqrt{65^2 - 33^2} = \sqrt{4225 - 1089} = \sqrt{3136} = 56$
\end{itemize}

Perimeter: $33 + 56 + 65 = 154$

\subsection*{C. Alternative Solution (Complex Numbers)}

Represent angles as arguments of complex numbers. The angle whose tangent is $\frac{a}{b}$ is the argument of $b + ai$.

For the three triangles:
\begin{itemize}
    \item $3-4-5$: angle with $\arg(4 + 3i)$
    \item $5-12-13$: angle with $\arg(12 + 5i)$
    \item Unknown: angle with $\arg(b + ai)$ where $a:b$ are the legs
\end{itemize}

Since the angles sum to $90^{\circ} = \frac{\pi}{2}$, their arguments sum to $\frac{\pi}{2}$, which is $\arg(i)$.

Therefore:
$$(4 + 3i)(12 + 5i)(b + ai) = ni$$
for some positive real number $n$.

Expand the left side:
\begin{align*}
    (4 + 3i)(12 + 5i) &= 48 + 20i + 36i + 15i^2 \\
    &= 48 + 56i - 15 \\
    &= 33 + 56i
\end{align*}

So:
$$(33 + 56i)(b + ai) = ni$$

Expand:
$$33b + 33ai + 56bi + 56ai^2 = ni$$
$$33b + 33ai + 56bi - 56a = ni$$
$$(33b - 56a) + (33a + 56b)i = ni$$

For this to be a pure imaginary number (real part = 0):
$$33b - 56a = 0$$
$$33b = 56a$$
$$\frac{a}{b} = \frac{33}{56}$$

Since we need a primitive triple, $a = 33$ and $b = 56$.

The hypotenuse is $\sqrt{33^2 + 56^2} = 65$.

Perimeter: $33 + 56 + 65 = 154$

\end{solutionbox}

\begin{warningbox}
\textbf{Common Pitfalls to Avoid}:

\begin{enumerate}
    \item \textbf{Wrong angle identification}: The smallest angle is opposite the \textit{smallest} leg, not the largest. For $3-4-5$, it's $\tan\alpha = \frac{3}{4}$, not $\frac{4}{3}$.
    
    \item \textbf{Tangent addition arithmetic errors}: 
    \begin{itemize}
        \item $\frac{3}{4} + \frac{5}{12} = \frac{14}{12}$ (not $\frac{8}{12}$)
        \item $1 - \frac{15}{48} = \frac{33}{48}$ (not $\frac{43}{48}$)
    \end{itemize}
    
    \item \textbf{Forgetting the undefined tangent}: Don't try to compute $\tan(90^{\circ})$ directly. Use the fact that it's undefined (denominator = 0).
    
    \item \textbf{Wrong tangent interpretation}: $\tan\theta = \frac{33}{56}$ means legs are 33 and 56, not that 33 is the hypotenuse.
    
    \item \textbf{Non-primitive triple confusion}: Always check $\gcd$ to ensure you have the primitive triple. Here $\gcd(33, 56) = 1$ \checkmark
    
    \item \textbf{Pythagorean calculation errors}:
    \begin{itemize}
        \item $33^2 = 1089$ (not 1099)
        \item $56^2 = 3136$ (not 3146)
        \item $\sqrt{4225} = 65$ (not 64 or 66)
    \end{itemize}
    
    \item \textbf{Perimeter addition error}: $33 + 56 + 65 = 154$ (not 144 or 164)
    
    \item \textbf{Complementary angle confusion}: If using $\theta = 90^{\circ} - (\alpha + \beta)$, remember $\sin\theta = \cos(\alpha + \beta)$, not $\sin(\alpha + \beta)$.
    
    \item \textbf{Answer choice misread}: Make sure to find the perimeter (sum of all three sides), not just the hypotenuse or sum of legs.
    
    \item \textbf{Complex number sign errors}: When expanding $(33 + 56i)(b + ai)$, remember $i^2 = -1$.
\end{enumerate}
\end{warningbox}

\begin{tipbox}
\textbf{Competition Tips}:

\begin{enumerate}
    \item \textbf{Recognize the "undefined tangent" setup}: When angles sum to $90^{\circ}$, $\tan(90^{\circ})$ is undefined. This means the denominator in tangent addition must be zero - direct path to the answer!
    
    \item \textbf{Tangent vs. Sine/Cosine choice}: 
    \begin{itemize}
        \item Tangent is faster (fewer calculations)
        \item Sine/cosine avoids the "undefined" concept but requires more arithmetic
        \item Choose based on what you're comfortable with
    \end{itemize}
    
    \item \textbf{Primitive Pythagorean triple shortcut}: Once you have $\tan\theta = \frac{33}{56}$ and verify $\gcd(33, 56) = 1$, the triple is automatically $33-56-65$. No need to question if there's scaling.
    
    \item \textbf{Quick $\gcd$ check}:
    \begin{itemize}
        \item $33 = 3 \cdot 11$ (odd factors)
        \item $56 = 8 \cdot 7$ (even, power of 2 times 7)
        \item No common factors $\Rightarrow$ $\gcd = 1$
    \end{itemize}
    
    \item \textbf{Hypotenuse calculation shortcut}: $\sqrt{33^2 + 56^2} = \sqrt{(33)(33) + (56)(56)}$. If you recognize that $33 = 3 \cdot 11$ and $56 = 8 \cdot 7$, you might recall the $33-56-65$ triple. Otherwise, calculate $\sqrt{4225} = 65$.
    
    \item \textbf{Answer choice reality check}: 
    \begin{itemize}
        \item (A) 40 is too small (less than $3 + 4 + 5 = 12$ or $5 + 12 + 13 = 30$)
        \item (B) 126 and (C) 154 are reasonable
        \item (D) 176 and (E) 208 seem large but possible
        \item Once you get legs 33 and 56, only (C) 154 works
    \end{itemize}
    
    \item \textbf{Complex number elegance}: If you're comfortable with complex numbers, this approach is very clean:
    $$(4 + 3i)(12 + 5i)(b + ai) = ni$$
    Real part = 0 immediately gives the answer!
    
    \item \textbf{Memorize common Pythagorean triples}: Besides $3-4-5$ and $5-12-13$, know:
    \begin{itemize}
        \item $8-15-17$
        \item $7-24-25$
        \item $20-21-29$
        \item $9-40-41$
        \item $\mathbf{33-56-65}$ (this problem!)
    \end{itemize}
    
    \item \textbf{Verification strategy}: After getting 154, quickly verify:
    \begin{itemize}
        \item $33 + 56 = 89$
        \item $89 + 65 = 154$ \checkmark
        \item Check $33^2 + 56^2 = 65^2$: $(33)(33) = 1089$, $(56)(56) = 3136$, sum $= 4225 = (65)^2$ \checkmark
    \end{itemize}
    
    \item \textbf{Time management}: This is P21, budget 4-6 minutes. Tangent addition is fastest (3-4 min). Complex numbers is elegant but requires comfort with the method (4-5 min). Sine/cosine is reliable but slower (5-6 min).
    
    \item \textbf{If stuck on tangent addition}: Switch to complementary angles. $\theta = 90^{\circ} - (\alpha + \beta)$ means $\cos\theta = \sin(\alpha + \beta)$. This avoids the "undefined" concept entirely.
\end{enumerate}
\end{tipbox}

\begin{competitionbox}
\subsection*{A. Time Management}
\begin{itemize}[leftmargin=*]
    \item \textbf{Estimated Time}: 4-6 minutes total
    \begin{itemize}
        \item Setup and angle identification: 30 seconds
        \item Tangent addition calculation: 2 minutes
        \item Solve for $\tan\theta$: 1 minute
        \item Find Pythagorean triple: 1 minute
        \item Verification: 30-60 seconds
    \end{itemize}
    
    \item \textbf{Time Checkpoints}: 
    \begin{itemize}
        \item At 1 min: Should have $\tan\alpha = \frac{3}{4}$ and $\tan\beta = \frac{5}{12}$
        \item At 3 min: Should have $\tan(\alpha + \beta) = \frac{56}{33}$
        \item At 4 min: Should have $\tan\theta = \frac{33}{56}$
        \item At 5 min: Should have triple $33-56-65$ and answer 154
    \end{itemize}
    
    \item \textbf{Priority Level}: \textbf{High} - This is P21 but very doable with standard techniques. If you know tangent addition, attempt this before harder problems.
\end{itemize}

\subsection*{B. Risk Assessment}
\begin{itemize}[leftmargin=*]
    \item \textbf{Confidence Level}: 90-95\% after verification
    \begin{itemize}
        \item Standard technique (tangent addition)
        \item Multiple solution paths available
        \item Pythagorean triple is verifiable
        \item Answer matches a choice
    \end{itemize}
    
    \item \textbf{Common Error Sources}:
    \begin{itemize}
        \item Arithmetic errors in tangent addition (fractions)
        \item Wrong angle identification (using wrong legs)
        \item Pythagorean theorem calculation errors
        \item Final perimeter addition error
        \item Not recognizing the triple as primitive
    \end{itemize}
    
    \item \textbf{Abandonment Criteria}: 
    \begin{itemize}
        \item If tangent addition is taking > 3 minutes $\Rightarrow$ try sine/cosine method
        \item If completely stuck after 5 minutes $\Rightarrow$ flag and move on
        \item If arithmetic is very messy $\Rightarrow$ check for calculation errors
    \end{itemize}
    
    \item \textbf{Flag for Review}: Yes if time permits - verify Pythagorean triple one more time
\end{itemize}

\subsection*{C. Strategic Context}
\begin{itemize}[leftmargin=*]
    \item \textbf{Problem Positioning}: P21 is accessible for students comfortable with trigonometry. The "undefined tangent" insight is key.
    
    \item \textbf{Score Impact}: Worth 6 points. High ROI if you recognize the tangent addition approach.
    
    \item \textbf{Time Budget Implications}: 4-6 minutes is reasonable. More efficient than typical P21 problems.
    
    \item \textbf{Technique Practice}: This problem tests:
    \begin{itemize}
        \item Right triangle angle relationships
        \item Tangent addition formula
        \item Undefined tangent condition
        \item Pythagorean triples
        \item Careful arithmetic with fractions
    \end{itemize}
\end{itemize}
\end{competitionbox}

\section*{Final Answer}

Using the tangent addition formula and the condition that $\tan(90^{\circ})$ is undefined:

\begin{itemize}
    \item $\tan\alpha = \frac{3}{4}$ (from $3-4-5$ triangle)
    \item $\tan\beta = \frac{5}{12}$ (from $5-12-13$ triangle)
    \item $\tan(\alpha + \beta) = \frac{56}{33}$
    \item $\tan\theta = \frac{33}{56}$ (from undefined tangent condition)
    \item Third triangle: $33-56-65$
    \item Perimeter: $33 + 56 + 65 = 154$
\end{itemize}

$$\boxed{\textbf{(C) } 154}$$

\section*{Confidence Assessment}
\begin{itemize}[leftmargin=*]
    \item \textbf{Confidence Level}: 95\%
    \begin{itemize}
        \item Tangent addition formula correctly applied
        \item Undefined tangent condition properly used
        \item Arithmetic verified at each step
        \item Pythagorean triple: $33^2 + 56^2 = 1089 + 3136 = 4225 = 65^2$ \checkmark
        \item Primitive check: $\gcd(33, 56) = 1$ \checkmark
        \item Perimeter: $33 + 56 + 65 = 154$ \checkmark
        \item Answer matches choice (C)
        \item Multiple solution methods confirm the same answer
    \end{itemize}
    
    \item \textbf{Verification Applied}: Level 4-5 (Standard to Thorough)
    \begin{itemize}
        \item Verified tangent addition arithmetic
        \item Checked Pythagorean triple calculation
        \item Confirmed primitive triple property
        \item Verified perimeter sum
        \item Cross-checked with alternative method (sine/cosine)
        \item All intermediate steps are correct
    \end{itemize}
    
    \item \textbf{Flag for Review}: No - high confidence, fully verified
    
    \item \textbf{Time Spent}: Approximately 5 minutes (within target range)
    \begin{itemize}
        \item 30 sec: Problem understanding and angle identification
        \item 2 min: Tangent addition calculation
        \item 1 min: Solving for third angle
        \item 1 min: Finding Pythagorean triple
        \item 30 sec: Final verification
    \end{itemize}
\end{itemize}

\section*{Learning Points}

\textbf{Key Takeaways from This Problem}:

\begin{enumerate}
    \item \textbf{Smallest Angle in Right Triangles}: Always opposite the smallest leg. This determines which tangent ratio to use.
    
    \item \textbf{Tangent Addition Formula}: Essential for angle sum problems:
    $$\tan(\alpha + \beta) = \frac{\tan\alpha + \tan\beta}{1 - \tan\alpha \tan\beta}$$
    
    \item \textbf{Undefined Tangent Technique}: When angles sum to $90^{\circ}$, $\tan(90^{\circ})$ is undefined. Set the denominator of the tangent addition formula to zero - this is faster than trying to compute the tangent directly!
    
    \item \textbf{Complementary Angle Relationships}: 
    \begin{itemize}
        \item $\sin(90^{\circ} - x) = \cos x$
        \item $\cos(90^{\circ} - x) = \sin x$
        \item Useful for alternative solution approaches
    \end{itemize}
    
    \item \textbf{Primitive Pythagorean Triples}: Once you have a ratio like $\frac{33}{56}$ with $\gcd = 1$, the triple is immediate. No scaling needed.
    
    \item \textbf{Complex Numbers for Angles}: The argument (angle) of $b + ai$ is $\arctan(\frac{a}{b})$. Product of complex numbers adds their arguments. Very elegant for angle sum problems!
    
    \item \textbf{Common Pythagorean Triples}: Memorize beyond just $3-4-5$ and $5-12-13$. The triple $33-56-65$ appears in competition problems.
    
    \item \textbf{Multiple Solution Paths}: This problem has at least 5 different elegant solutions. Knowing multiple methods gives you flexibility in competition.
\end{enumerate}

\vspace{0.5cm}

\textbf{Related Problems to Practice}:
\begin{itemize}
    \item Other AMC12 problems involving angle sums in right triangles
    \item Tangent addition and complementary angle problems
    \item Pythagorean triple identification
    \item Complex number arguments and angle representations
\end{itemize}

\vspace{0.5cm}

\textbf{Essential Formulas}:

\textbf{Tangent Addition Formula}:
$$\tan(\alpha + \beta) = \frac{\tan\alpha + \tan\beta}{1 - \tan\alpha \tan\beta}$$

\textbf{Complementary Angle Identities}:
\begin{align*}
    \sin(90^{\circ} - x) &= \cos x \\
    \cos(90^{\circ} - x) &= \sin x \\
    \tan(90^{\circ} - x) &= \cot x = \frac{1}{\tan x}
\end{align*}

\textbf{Sine/Cosine Addition}:
\begin{align*}
    \sin(\alpha + \beta) &= \sin\alpha \cos\beta + \cos\alpha \sin\beta \\
    \cos(\alpha + \beta) &= \cos\alpha \cos\beta - \sin\alpha \sin\beta
\end{align*}

\textbf{Common Primitive Pythagorean Triples}:
\begin{itemize}
    \item $(3, 4, 5)$
    \item $(5, 12, 13)$
    \item $(8, 15, 17)$
    \item $(7, 24, 25)$
    \item $(20, 21, 29)$
    \item $(9, 40, 41)$
    \item \textbf{$(33, 56, 65)$} (this problem!)
\end{itemize}

\end{document}

